\section{Limitations}
\label{sec:limitations}

\begin{enumerate}
\item \emph{The pooling assumption is not tested.} Testing it model by model needs labels drawn
  by model within the concordant cells, which the 400-item audit, stratified otherwise, does not
  provide, and the reference tolerance is a judgement, not calibrated against any external study
  (Section~\ref{sec:setup}). If the
  assumption fails at that tolerance, the certified counts are too high, which is why every result
  is reported as a curve over $\varepsilon$, with a breakdown value for each ordering and the
  no-pooling analysis beside it.
\item \emph{The reader verdict inherits the benchmark's reference solutions.} Every verifier and
  the reader verdict grade against the same reference solutions: the harness derives its gold
  answer from each solution's last boxed expression, and re-deriving it agrees on
  \GoldConsistentN{} of \NumItems{} items; a normalisation rule that damages \GoldFtItems{} of those
  answers changes no audited verdict (Appendix~\ref{sec:sensitivity}). An error in a reference solution therefore affects our
  results and the leaderboard's in the same way, and no analysis inside the pipeline can detect it.
  An audit of a sample of reference solutions would.
\item \emph{The verifiers share machinery.} A response that defeats the shared symbolic library
  could be miscounted by all three verifiers and land in a concordant cell, which would hide
  disagreement rather than create it. xFinder, which extracts differently, moved the count by one.
\item \emph{The scope is one benchmark on one archived board.} The roster of \NumModels{} models
  is a convenience sample that sits low on the board and moved less than the board under the
  change of verifier, so its counts are unlikely to overstate the problem but do not estimate it;
  the secondary analysis also treats a curated benchmark as a random sample.
\item \emph{The labels come mainly from one labeller, the author.} On overlaps of \LabTwoN{} and
  \LabThreeN{} items, two others agree with the reader verdicts in \LabTwoPct\% and \LabThreePct\%
  of cases, and the labels carry no record of the time spent on each item (Appendix~\ref{app:data}).
\end{enumerate}

\section{Conclusion}
\label{sec:conclusion}

A leaderboard names the model but not the verifier, and on this roster the choice of verifier
moves the ranking without the published score showing it. A small human audit can settle part of
the ranking, but only under an explicit assumption about how models differ where the verifiers
agree. Without that assumption this audit certifies \CertFixedNoPooling{} of the \NumOrderings{}
orderings on these problems; at the reference tolerance it certifies \CertFixedPrereg{}, all at once
and with a finite-sample guarantee. The honest output is the curve between the two, with a breakdown
value for every ordering.

Most disagreements here were responses in which a verifier found no answer. Publishing, beside
each score, the verifier, its version and that count, which the pipeline computes and then
discards, would let a reader bound how far a score could rise if every unread response were correct.
It would not settle the disagreements that remain after extraction, which only a reader can. The
assumption this paper rests on can be tested directly with labels drawn model by model in the
concordant cells, and that is the natural next step.
